\documentclass[a4paper,11pt]{article}

% ============================================================
%  Le Collé — Séance 1 : Déterminants, courbes paramétrées
%  © 2026 Ahmed Amine Ouqlani
% ============================================================

% ---------- Polices ----------
\usepackage[utf8]{inputenc}
\usepackage[T1]{fontenc}
\usepackage[sc]{mathpazo}          % texte et mathématiques : Palatino
\usepackage[scaled=0.88]{helvet}   % titres et repères : linéale
\linespread{1.07}
\usepackage{amsmath,amssymb,amsthm,mathrsfs,stmaryrd}
\usepackage{microtype}

% ---------- Mise en page ----------
\usepackage[a4paper,top=2.3cm,bottom=2.3cm,left=2.4cm,right=2.4cm,
            headheight=14pt,headsep=0.8cm,footskip=1.1cm]{geometry}
\setlength{\parindent}{0pt}
\setlength{\parskip}{0.45em plus 0.15em minus 0.1em}
\raggedbottom
\emergencystretch=1.5em
\clubpenalty=10000 \widowpenalty=10000 \displaywidowpenalty=10000 \brokenpenalty=10000

% ---------- Couleurs : un bleu de nuit, un ocre, des gris ----------
\usepackage{xcolor}
\definecolor{Ink}{HTML}{1C2533}
\definecolor{Bleu}{HTML}{1F3A5F}
\definecolor{BleuMoyen}{HTML}{4A6A8F}
\definecolor{BleuClair}{HTML}{9DB2C9}
\definecolor{Fond}{HTML}{F2F4F7}
\definecolor{Ocre}{HTML}{B07A26}
\definecolor{Brique}{HTML}{A23B2A}
\definecolor{Gris}{HTML}{5E6672}
\definecolor{Filet}{HTML}{C5CCD6}

% ---------- Packages ----------
\usepackage{enumitem}
\setlist{itemsep=0.25em,topsep=0.35em,parsep=0pt}
\setlist[itemize]{label=\textcolor{BleuMoyen}{\small\textbullet},leftmargin=1.4em}
\setlist[enumerate]{leftmargin=1.8em,label=\textsf{\small\color{Bleu}\arabic*.}}
\usepackage{booktabs,array,multicol}
\usepackage{needspace,float}
\usepackage[font=small,labelfont={sf,color=Bleu},labelsep=period,skip=6pt]{caption}
\usepackage{tikz}
\usetikzlibrary{arrows.meta,calc,decorations.markings}
\usepackage{pgfplots}
\pgfplotsset{compat=1.17}
\usepackage{tkz-tab}
\usepackage[most]{tcolorbox}
\usepackage{titlesec,titletoc}
\usepackage{fancyhdr,lastpage,etoolbox}
\usepackage[colorlinks=true,linkcolor=Bleu,urlcolor=Bleu,citecolor=Bleu,
            pdftitle={Le Collé — TD Équations différentielles (corrections)},pdfauthor={Ahmed Amine Ouqlani}]{hyperref}

\numberwithin{equation}{section}
\numberwithin{figure}{section}
\numberwithin{table}{section}

% ---------- Titres ----------
% Chapitres (sections) : grand numéro, titre en linéale, filet ocre
\titleformat{\section}[display]
  {\normalfont\sffamily\color{Ink}}
  {{\small\color{Gris}\textls[200]{\MakeUppercase{Chapitre}}}\\[-0.2em]%
   {\fontsize{64}{64}\selectfont\color{BleuClair}\thesection}}
  {0.4em}
  {\fontsize{25}{30}\selectfont\bfseries}
  [\vspace{0.5em}{\color{Ocre}\rule{3.2cm}{1.4pt}}]
\titlespacing*{\section}{0pt}{0pt}{1.6em}
\titleformat{name=\section,numberless}[display]
  {\normalfont\sffamily\color{Ink}}{}{0pt}{\fontsize{25}{30}\selectfont\bfseries}
  [\vspace{0.5em}{\color{Ocre}\rule{3.2cm}{1.4pt}}]
\titlespacing*{name=\section,numberless}{0pt}{0pt}{1.6em}

\titleformat{\subsection}
  {\needspace{5\baselineskip}\normalfont\sffamily\Large\bfseries\color{Ink}}
  {\color{Bleu}\thesubsection}{0.8em}{}
\titlespacing*{\subsection}{0pt}{1.8em}{0.5em}
\titleformat{\subsubsection}
  {\needspace{5\baselineskip}\normalfont\sffamily\large\bfseries\color{Ink}}
  {\color{Bleu}\thesubsubsection}{0.7em}{}
\titlespacing*{\subsubsection}{0pt}{1.6em}{0.4em}
\titleformat{\paragraph}[runin]{\normalfont\sffamily\bfseries\color{Ink}}{}{0pt}{}[.]
\titlespacing*{\paragraph}{0pt}{0.9em}{0.6em}
\setcounter{secnumdepth}{3}
\setcounter{tocdepth}{2}

% ---------- Table des matières ----------
\titlecontents{section}[0pt]{\addvspace{0.7em}\sffamily\bfseries\color{Ink}}
  {\makebox[2.2em][l]{\color{Bleu}\thecontentslabel}}{}{\hfill\contentspage}[\addvspace{0.2em}]
\titlecontents{subsection}[2.2em]{\small\linespread{1}\selectfont}
  {\makebox[2.6em][l]{\color{Gris}\thecontentslabel}}{}{\titlerule*[0.7pc]{.}\contentspage}

% ---------- En-têtes et pieds de page ----------
\pagestyle{fancy}
\fancyhf{}
\renewcommand{\sectionmark}[1]{\markboth{\thesection\quad#1}{}}
\fancyhead[L]{\footnotesize\sffamily\color{Gris}\leftmark}
\fancyhead[R]{\footnotesize\sffamily\scshape\color{Gris}Le Collé}
\fancyfoot[C]{\small\sffamily\color{Ink}\thepage}
\fancyfoot[L]{\scriptsize\sffamily\hyperlink{toc}{\color{Filet}sommaire}}
\renewcommand{\headrulewidth}{0.4pt}
\renewcommand{\headrule}{{\color{Filet}\hrule height\headrulewidth width\headwidth}}
\renewcommand{\footrulewidth}{0pt}
\fancypagestyle{plain}{\fancyhf{}\fancyfoot[C]{\small\sffamily\color{Ink}\thepage}\renewcommand{\headrulewidth}{0pt}}

% ---------- Énoncés ----------
\newtheoremstyle{enonce}{0pt}{0pt}{}{}{\sffamily\bfseries\color{Bleu}}{.}{0.5em}
  {\thmname{#1}\thmnumber{~#2}\thmnote{\ \normalfont\sffamily\color{Ink}(#3)}}
\newtheoremstyle{exemple}{0.9em}{0.9em}{}{}{\sffamily\bfseries\color{BleuMoyen}}{.}{0.5em}
  {\thmname{#1}\thmnumber{~#2}\thmnote{\ \normalfont\sffamily\color{Ink}(#3)}}
\newtheoremstyle{remarque}{0.7em}{0.7em}{}{}{\sffamily\itshape}{.}{0.5em}{\thmname{#1}\thmnote{\ (#3)}}
\theoremstyle{enonce}
\newtheorem{definition}{Définition}[section]
\newtheorem{theoreme}[definition]{Théorème}
\newtheorem{proposition}[definition]{Proposition}
\newtheorem{corollaire}[definition]{Corollaire}
\theoremstyle{exemple}
\newtheorem{exemple}[definition]{Exemple}
\theoremstyle{remarque}
\newtheorem*{remarque}{Remarque}

\tcbset{
  defi/.style={enhanced,breakable,colback=white,colframe=white,boxrule=0pt,arc=0pt,
     borderline west={1.2pt}{0pt}{Bleu},left=11pt,right=4pt,top=3pt,bottom=3pt,
     before skip=1.1em,after skip=1.1em},
  thm/.style={enhanced,breakable,colback=Fond,colframe=Fond,boxrule=0pt,arc=0pt,
     borderline west={2.4pt}{0pt}{Bleu},left=11pt,right=8pt,top=6pt,bottom=6pt,
     before skip=1.2em,after skip=1.2em}}
\tcolorboxenvironment{definition}{defi}
\tcolorboxenvironment{theoreme}{thm}
\tcolorboxenvironment{proposition}{thm}
\tcolorboxenvironment{corollaire}{thm}

\renewcommand{\proofname}{Démonstration}
\renewcommand{\qedsymbol}{\textcolor{BleuClair}{$\blacksquare$}}
\makeatletter
\renewenvironment{proof}[1][\proofname]{\par\pushQED{\qed}\normalfont
  \topsep0.5em\relax\trivlist\item[\hskip\labelsep\sffamily\itshape\color{Gris}#1.]\ignorespaces}
  {\popQED\endtrivlist\@endpefalse}
\makeatother

% ---------- Méthode : filets haut et bas, titre en bandeau ----------
\makeatletter\newcommand{\methlabelfix}{\protected@edef\@currentlabel{\thetcbcounter}}\makeatother
\newtcolorbox[auto counter,number within=section]{methode}[2][]{%
  enhanced,unbreakable,colback=white,colframe=white,boxrule=0pt,arc=0pt,
  borderline north={0.9pt}{0pt}{Bleu},borderline south={0.9pt}{0pt}{Bleu},
  left=4pt,right=4pt,top=8pt,bottom=8pt,before skip=1.4em,after skip=1.4em,
  before upper={\methlabelfix\phantomsection\addcontentsline{met}{metentry}{\protect\makebox[2.6em][l]{\thetcbcounter}#2}%
     {\sffamily\bfseries\small\color{white}\colorbox{Bleu}{\strut\ MÉTHODE~\thetcbcounter\ }}%
     \ \ {\sffamily\bfseries\color{Bleu}#2}\par\smallskip},#1}

% ---------- Repères pédagogiques : un filet de couleur, un mot-clé ----------
\newcommand{\definirbloc}[3]{% nom, couleur, libellé
  \newtcolorbox{#1}[1][]{enhanced,breakable,blanker,
    borderline west={1.4pt}{0pt}{#2},left=10pt,right=2pt,top=1pt,bottom=1pt,
    before skip=1em,after skip=1em,
    before upper={{\sffamily\bfseries\footnotesize\color{#2}\textls[120]{\MakeUppercase{#3}}}\quad},##1}}
\definirbloc{reflexe}{Bleu}{Réflexe}
\definirbloc{astuce}{Ocre}{Astuce}
\definirbloc{attention}{Brique}{Attention}
\definirbloc{enpratique}{BleuMoyen}{En pratique}
\definirbloc{plusloin}{Gris}{Pour aller plus loin}
\definirbloc{strategie}{Bleu}{Stratégie}
\definirbloc{lecture}{BleuMoyen}{Lire le schéma}
% compatibilité avec les parties reprises
\newenvironment{garde}{\begin{attention}}{\end{attention}}
\newenvironment{complementHP}{\begin{plusloin}}{\end{plusloin}}
\newenvironment{interpretation}{\begin{enpratique}}{\end{enpratique}}

% ---------- Exercices, corrigés ----------
\newcounter{exo}[section]
\renewcommand{\theexo}{\thesection.\arabic{exo}}
\newcommand{\niveau}[1]{{\color{Ocre}\footnotesize
  \ifcase#1\or$\bullet\circ\circ\circ\circ$\or$\bullet\bullet\circ\circ\circ$\or$\bullet\bullet\bullet\circ\circ$%
  \or$\bullet\bullet\bullet\bullet\circ$\or$\bullet\bullet\bullet\bullet\bullet$\fi}}
\newcommand{\serie}[1]{\par\needspace{8\baselineskip}\vspace{1.5em}%
  \noindent{\sffamily\bfseries\color{Bleu}#1}\par\nobreak\vspace{-0.4em}%
  \noindent{\color{Filet}\rule{\linewidth}{0.4pt}}\par\nobreak\vspace{0.1em}}
\newenvironment{exo}[3]{%
  \par\needspace{5\baselineskip}\vspace{0.75em}\refstepcounter{exo}\label{e:#1}%
  \def\exolab{#1}%
  \noindent{\sffamily\bfseries\color{Bleu}Exercice~\theexo}\quad\niveau{#2}%
  \hfill{\footnotesize\sffamily\color{Gris}#3}\par\nobreak\vspace{0.15em}}%
  {\par\nopagebreak\noindent\hfill{\scriptsize\sffamily\hyperref[c:\exolab]{\color{Gris}corrigé p.~\pageref{c:\exolab}}}\par}
\newenvironment{corrige}[1]{%
  \par\needspace{6\baselineskip}\vspace{1em}\phantomsection\label{c:#1}%
  \noindent{\sffamily\bfseries\color{Bleu}Exercice~\ref{e:#1}}\quad
  {\scriptsize\sffamily\hyperref[e:#1]{\color{Gris}énoncé p.~\pageref{e:#1}}}\par\nobreak\vspace{0.2em}}%
  {\hfill\textcolor{BleuClair}{$\square$}\par}
\newcommand{\etape}[1]{\par\smallskip\noindent{\sffamily\bfseries\small\color{BleuMoyen}#1.}\ }
\newenvironment{solution}{\par\noindent{\sffamily\bfseries\small\color{BleuMoyen}Solution.}\ }{\hfill\textcolor{BleuClair}{$\square$}\par}
\newcommand{\chercher}{\par\smallskip\noindent{\color{Filet}\rule[0.5ex]{1.2cm}{0.4pt}}\
  {\footnotesize\sffamily\itshape\color{Gris}laisser chercher l'étudiant avant de poursuivre}\par\smallskip}

% ---------- Guide de séance ----------

\newcommand{\marge}[1]{\leavevmode\marginpar{\raggedleft\scriptsize\sffamily\color{Gris}\textls[100]{\MakeUppercase{#1}}}\ignorespaces}
\newcommand{\Etape}[4]{%
  \subsection[Étape #1 --- #3 {\color{Gris}(#2)}]{#3}\label{etape:#1}%
  \noindent{\small\sffamily\color{Gris}{\color{Ocre}\bfseries Étape #1}\quad·\quad #2\quad·\quad #4}\par\medskip}
\newcommand{\question}[1]{\par\smallskip\noindent{\sffamily\bfseries\small\color{BleuMoyen}Question.}\ \textit{#1}\par\smallskip}
\newcommand{\rappel}[1]{\par\smallskip\noindent\phantomsection\addcontentsline{rap}{rapentry}{#1}%
  {\sffamily\bfseries\scriptsize\colorbox{Fond}{\color{Gris}\strut RAPPEL}}\ {\sffamily\small\color{Gris}#1.}\ }
\newcommand{\pourquoi}{\par\noindent{\sffamily\bfseries\small\color{Ocre}Pourquoi ?}\ }

% Ouverture de chapitre : \ouverture{objectifs}{prérequis}{durée}
\newcommand{\ouverture}[3]{%
  \begin{tcolorbox}[enhanced,colback=Fond,colframe=Fond,arc=0pt,boxrule=0pt,left=12pt,right=12pt,top=10pt,bottom=10pt,
     before skip=1.2em,after skip=1.6em]
  \begin{minipage}[t]{0.56\linewidth}
  {\sffamily\bfseries\small\color{Bleu}OBJECTIFS}\par\vspace{0.3em}\small #1
  \end{minipage}\hfill
  \begin{minipage}[t]{0.38\linewidth}
  {\sffamily\bfseries\small\color{Bleu}PRÉREQUIS}\par\vspace{0.3em}\small #2
  \ifx\relax#3\relax\else\par\vspace{0.6em}{\sffamily\bfseries\small\color{Bleu}DURÉE}\par\vspace{0.2em}\small #3\fi
  \end{minipage}
  \end{tcolorbox}}
% Fin de chapitre
\newenvironment{essentiel}{\par\needspace{12\baselineskip}\vspace{1.5em}%
  \begin{tcolorbox}[enhanced,unbreakable,colback=white,colframe=white,boxrule=0pt,arc=0pt,
   borderline north={1.4pt}{0pt}{Ocre},left=2pt,right=2pt,top=8pt,bottom=4pt]
  {\sffamily\bfseries\color{Ocre}L'ESSENTIEL DU CHAPITRE}\par\medskip}{\end{tcolorbox}}

% ---------- Notations ----------
\newcommand{\K}{\mathbb{K}}
\newcommand{\R}{\mathbb{R}}
\newcommand{\C}{\mathbb{C}}
\newcommand{\N}{\mathbb{N}}
\newcommand{\Z}{\mathbb{Z}}
\newcommand{\Mn}[1]{\mathcal{M}_{n}(#1)}
\newcommand{\GLn}[1]{\mathrm{GL}_n(#1)}
\newcommand{\LE}{\mathcal{L}(E)}
\newcommand{\rg}{\operatorname{rg}}
\newcommand{\vect}{\operatorname{Vect}}
\newcommand{\Mat}{\operatorname{Mat}}
\newcommand{\Id}{\operatorname{Id}}
\newcommand{\sh}{\operatorname{sh}}
\newcommand{\ch}{\operatorname{ch}}
\renewcommand{\th}{\operatorname{th}}
\DeclareMathOperator{\argsh}{argsh}
\newcommand{\vv}[1]{\overrightarrow{#1}}
\newcommand{\ve}[1]{\vec{#1}}
\newcommand{\ient}[2]{\llbracket #1,#2\rrbracket}
\newcommand{\dt}[1]{\left|\begin{matrix}#1\end{matrix}\right|}
\newcommand{\dis}{\displaystyle}

% ---------- Figures ----------
\pgfplotsset{trig format plots=rad}
\pgfplotsset{courbe/.style={axis lines=middle,axis equal image,
  axis line style={Gris,-{Stealth[length=5pt]}},
  xlabel={$x$},ylabel={$y$},ticklabel style={font=\scriptsize\sffamily,color=Gris},
  every axis x label/.style={at={(ticklabel* cs:1)},anchor=west,font=\small},
  every axis y label/.style={at={(ticklabel* cs:1)},anchor=south,font=\small},
  samples=300,smooth,width=0.6\linewidth}}
\tikzset{arcetude/.style={Bleu,line width=1.25pt},
         reste/.style={BleuClair,line width=1pt},
         tgte/.style={Ocre,line width=0.9pt},
         asy/.style={Gris,dashed,line width=0.6pt},
         pt/.style={circle,fill=Ink,inner sep=1.2pt}}

% ---------- Navigation ----------
\makeatletter
\newcommand*\l@rapentry{\@dottedtocline{1}{0em}{0em}}
\newcommand*\l@metentry{\@dottedtocline{1}{0em}{0em}}
\newcommand{\listemethodes}{\@starttoc{met}}
\newcommand{\listerappels}{\@starttoc{rap}}
\makeatother
% onglets de chapitre sur le bord droit
\usepackage{eso-pic}
\newif\ifonglets \ongletsfalse
\newcommand{\dessineronglet}{%
  \edef\ongcol{\ifnum\value{section}=3 Ocre\else Bleu\fi}%
  \edef\ongnum{\the\value{section}}%
  \pgfmathsetmacro{\ongy}{3.2+2.3*(\ongnum-1)}%
  \begin{tikzpicture}[remember picture,overlay]
  \fill[\ongcol] ([xshift=-0.75cm,yshift=-\ongy cm]current page.north east)
     rectangle ([yshift=-\ongy cm-1.9cm]current page.north east);
  \node[text=white,font=\sffamily\bfseries\large] at ([xshift=-0.37cm,yshift=-\ongy cm-0.95cm]current page.north east)
     {\ongnum};
  \end{tikzpicture}}
\AddToShipoutPictureBG{\ifonglets\ifnum\value{section}>0 \dessineronglet\fi\fi}
% blocs de navigation et d'enseignement
\definirbloc{commencer}{Gris}{Si vous commencez ici}
\definirbloc{horsprogramme}{Brique}{Hors programme}
\newcommand{\HP}{\,{\sffamily\bfseries\scriptsize\textcolor{Brique}{[HP]}}}
\newcommand{\source}[1]{{\footnotesize\sffamily\color{Gris}\,(#1)}}
\newcommand{\coupdepouce}[1]{\par\vspace{-0.2em}\noindent{\footnotesize\sffamily\color{Gris}%
  \textbf{Coup de pouce} (à donner si l'étudiant bloque) : #1}\par}
\newcommand{\pointarret}[1]{\par\vspace{1em}\noindent{\color{Filet}\rule[0.5ex]{\linewidth}{0.4pt}}\par\nobreak
  \noindent{\footnotesize\sffamily\color{Ocre}\bfseries POINT D'ARRÊT POSSIBLE}\quad{\footnotesize\sffamily\color{Gris}#1}\par\vspace{0.6em}}
\newcommand{\demarche}[1]{\noindent{\sffamily\bfseries\small\color{BleuMoyen}Démarche.}\ \textit{#1}\par\smallskip}

% ---------- Corrigés : idées, astuces, pièges, à retenir ----------
\newcommand{\choix}[1]{\par\noindent{\sffamily\bfseries\small\color{BleuMoyen}Choix de la méthode.}\ #1\par\smallskip}
\newcommand{\gain}[1]{\par\noindent{\sffamily\bfseries\small\color{Ocre}Astuce.}\ #1\par}
\newcommand{\piege}[1]{\par\noindent{\sffamily\bfseries\small\color{Brique}Piège.}\ #1\par}
\newtcolorbox{retenirbox}{enhanced,breakable,colback=Fond,colframe=Fond,arc=0pt,boxrule=0pt,
  borderline west={2pt}{0pt}{Ocre},left=8pt,right=6pt,top=4pt,bottom=4pt,before skip=0.7em,after skip=0.5em}
\newcommand{\retenir}[1]{\begin{retenirbox}{\sffamily\bfseries\footnotesize\color{Ocre}À RETENIR}\quad #1\end{retenirbox}}
\newcommand{\TD}[1]{{\color{Ocre}\bfseries TD\,#1}\ ·\ }
\newcommand{\q}[1]{\par\smallskip\noindent{\sffamily\bfseries\color{Bleu}#1}\ }
\providecommand{\iddots}{\mathinner{\mkern1mu\raise1pt\hbox{.}\mkern2mu\raise4pt\hbox{.}\mkern2mu\raise7pt\hbox{.}\mkern1mu}}

\providecommand{\Ker}{\operatorname{Ker}}\renewcommand{\Im}{\operatorname{Im}}\providecommand{\tr}{\operatorname{tr}}\providecommand{\Sp}{\operatorname{Sp}}

% ================================================================ TD : mise en page propre au livret d'exercices
\fancyhf{}
\fancyhead[L]{\footnotesize\sffamily\color{Gris}TD · Équations différentielles · corrections}
\fancyhead[R]{\footnotesize\sffamily\scshape\color{Gris}Le Collé}
\fancyfoot[C]{\small\sffamily\color{Ink}\thepage\,/\,\pageref{LastPage}}
\fancypagestyle{plain}{\fancyhf{}\fancyfoot[C]{\small\sffamily\color{Ink}\thepage\,/\,\pageref{LastPage}}\renewcommand{\headrulewidth}{0pt}}
\makeatletter
\newcommand{\partie}[3]{% numéro, titre, description
  \par\needspace{10\baselineskip}\vspace{1.4em}\phantomsection\label{partie:#1}\pdfbookmark[0]{Partie #1}{partie.#1}%
  {\sffamily\small\color{Gris}\textls[200]{PARTIE #1}}\par\nobreak\vspace{0.1em}
  {\sffamily\bfseries\LARGE\color{Ink}#2}\par\nobreak\vspace{0.3em}
  {\color{Ocre}\rule{2.6cm}{1.2pt}}\par\nobreak\vspace{0.2em}
  {\small\color{Gris}#3}\par\nobreak\vspace{0.6em}}
\newcommand{\exercice}[7]{% numéro, difficulté, titre, origine, notions, identifiant, page du corrigé
  \par\needspace{8\baselineskip}\vspace{1.1em}\phantomsection
  \def\@currentlabel{#1}\label{e:#6}\pdfbookmark[1]{Exercice #1}{e.#6}%
  \noindent{\sffamily\bfseries\large\color{Bleu}Exercice~#1}\quad\niveau{#2}\hfill{\footnotesize\sffamily\color{Gris}#4}\par\nobreak
  \vspace{0.1em}\noindent{\sffamily\bfseries\color{Ink}#3}\par\nobreak
  \vspace{0.1em}\noindent{\footnotesize\sffamily\color{Gris}Notions : #5.\hfill\mbox{corrigé p.~#7}\par}\nobreak
  \vspace{0.35em}}
\newcommand{\renvoi}[1]{\par\nopagebreak\vspace{0.2em}\noindent\hfill{\scriptsize\sffamily\color{Gris}#1}\par}
\newcommand{\place}{\par\vspace{1.6cm plus 0.6cm minus 0.6cm}}
\newcommand{\correction}[4]{% numéro, titre, identifiant, renvoi
  \par\needspace{7\baselineskip}\vspace{1.3em}\phantomsection
  \def\@currentlabel{#1}\label{e:#3}\pdfbookmark[1]{Exercice #1}{c.#3}%
  {\color{Filet}\hrule height 0.4pt}\vspace{0.5em}
  \noindent{\sffamily\bfseries\large\color{Bleu}Exercice~#1}\quad{\sffamily\bfseries\color{Ink}#2}\hfill{\scriptsize\sffamily\color{Gris}#4}\par\nobreak
  \vspace{0.4em}}
\newcommand{\sommaireligne}[3]{\noindent\hbox to 2.2em{\hfill\sffamily\color{Bleu}#1.}\enspace\begin{minipage}[t]{0.8\linewidth}\raggedright#2\end{minipage}\hfill\pageref{e:#3}\par}
\makeatother
\setlength{\parskip}{0.5em plus 0.15em minus 0.1em}

\begin{document}
\thispagestyle{plain}
{\sffamily\small\color{Gris}\textls[150]{LE COLLÉ · MATHÉMATIQUES PT · CHAPITRE 4}}\par\vspace{0.2em}
{\sffamily\bfseries\fontsize{26}{30}\selectfont\color{Ink}TD --- Équations différentielles}\par\vspace{0.3em}
{\sffamily\Large\color{BleuMoyen}Corrections}\par\vspace{0.4em}
{\color{Ocre}\rule{3.2cm}{1.4pt}}\par\vspace{0.8em}
{\small Corrections détaillées des 17 exercices du TD, dans le même ordre et avec les mêmes numéros.
Chaque correction commence par le choix de la méthode, puis détaille les calculs ; les encadrés « À retenir » résument ce qui
resservira. Les renvois « méthode~1.3 du livret » désignent le cours.
\textbf{Chercher d'abord l'exercice} avant d'ouvrir sa correction.}

\partie{I}{Applications directes du cours}{}
\correction{1}{Premier ordre et solution constante}{edn-1}{énoncé : TD, p.~2}
Homogène : une primitive de $2t$ est $t^2$, d'où $y=Ke^{-t^2}$. La constante $\frac12$ est solution évidente ($0+2t\cdot\frac12=t$).
Solutions : $y=\frac12+Ke^{-t^2}$ ; $y(0)=0$ donne $K=-\frac12$, soit $y=\frac12\big(1-e^{-t^2}\big)$.
\gain{toujours tester une constante quand le second membre est proportionnel au coefficient de $y$.}
\par\hfill\textcolor{BleuClair}{$\square$}\par
\correction{2}{Solutions particulières au premier ordre}{ed-o1}{énoncé : TD, p.~2}
\q{(a)} Homogène : $Ke^{-t}$. Par $z'+z=e^{it}$, $z=\frac{e^{it}}{1+i}=\frac{(1-i)e^{it}}2$, d'où $y_p=\operatorname{Im}z=\frac{\sin t-\cos t}2$.
Solutions : $y=\frac{\sin t-\cos t}2+Ke^{-t}$.

\q{(b)} Homogène : $Ke^{2t}$. Le second membre est $\operatorname{Im}\big(e^{(1+2i)t}\big)$ et $1+2i\neq2$ : on cherche $z=\nu e^{(1+2i)t}$,
$\nu(1+2i-2)=1$, $\nu=\frac1{-1+2i}=\frac{-1-2i}5$. Donc $y_p=\operatorname{Im}z=-\frac{e^t}5\big(\sin2t+2\cos2t\big)$ et
$y=Ke^{2t}-\frac{e^t}5(\sin2t+2\cos2t)$.

\q{(c)} $m=-1=-a$ : résonance, on cherche $y=Q(t)e^{-t}$ avec $\deg Q=2$. En reportant : $Q'=t$, $Q=\frac{t^2}2$.
Solutions : $y=\left(\frac{t^2}2+K\right)e^{-t}$.
\piege{en (c), chercher $(\alpha t+\beta)e^{-t}$ échoue : le terme $\beta e^{-t}$ est solution de l'homogène.}
\par\hfill\textcolor{BleuClair}{$\square$}\par
\correction{3}{Normaliser avant de résoudre}{edn-2}{énoncé : TD, p.~2}
Forme normalisée : $y'+\frac1{2t}y=\frac{1+t}{2t}$. Une primitive de $\frac1{2t}$ est $\frac12\ln t$ : homogène $y=Ke^{-\frac12\ln t}=\frac K{\sqrt t}$.
On cherche une solution affine $y=a+bt$ : $2bt+a+bt=1+t$, d'où $a=1$, $b=\frac13$. Solutions : $y=1+\frac t3+\frac K{\sqrt t}$.
\par\hfill\textcolor{BleuClair}{$\square$}\par
\correction{4}{Trois équations homogènes}{ed-c1}{énoncé : TD, p.~2}
(a) $r^2-3r+2=(r-1)(r-2)$ : $y=Ae^t+Be^{2t}$. (b) $(r+1)^2$ : $y=(A+Bt)e^{-t}$.
(c) $r=-\frac12\pm i\frac{\sqrt3}2$ : $y=e^{-t/2}\left(A\cos\frac{\sqrt3}2t+B\sin\frac{\sqrt3}2t\right)$.
\par\hfill\textcolor{BleuClair}{$\square$}\par
\correction{5}{Deux seconds membres}{edn-16}{énoncé : TD, p.~3}
$S_H=\{A\cos t+B\sin t\}$ (racines $\pm i$). Par superposition : $y=t$ convient pour le second membre $t$, et $y=ce^t$ donne
$2c=1$ pour $e^t$. Solutions : $y=t+\frac12e^t+A\cos t+B\sin t$.
\par\hfill\textcolor{BleuClair}{$\square$}\par
\correction{6}{Un problème de Cauchy du premier ordre}{ed-cauchy1}{énoncé : TD, p.~3}
Sur $\left]-\frac\pi2,\frac\pi2\right[$, $\cos t\neq0$ : $y'+\tan(t)\,y=\frac1{\cos t}$. Une primitive de $\tan$ est $-\ln(\cos t)$, d'où
l'homogène $y=Ke^{\ln\cos t}=K\cos t$. Variation de la constante : $y=\lambda\cos t$ donne $\lambda'\cos t=\frac1{\cos t}$,
$\lambda=\tan t+K$. Solutions : $y=\sin t+K\cos t$ ; $y(0)=K=2$. La solution est $y=\sin t+2\cos t$.
\gain{une solution particulière évidente ($\sin t$) se vérifie en une ligne : $\cos t\cos t+\sin t\sin t=1$.}
\par\hfill\textcolor{BleuClair}{$\square$}\par
\correction{7}{Problème de Cauchy, racines complexes}{edn-4}{énoncé : TD, p.~3}
$r^2+2r+5=0$ : $r=-1\pm2i$. Donc $y=e^{-t}(A\cos2t+B\sin2t)$. $y(0)=A=1$ ; $y'(0)=-A+2B=1$, $B=1$.
$y=e^{-t}(\cos2t+\sin2t)=\sqrt2\,e^{-t}\cos\left(2t-\frac\pi4\right)$ : oscillations de pulsation $2$, amorties, qui tendent vers $0$.
\par\hfill\textcolor{BleuClair}{$\square$}\par
\correction{8}{Résonance et problème de Cauchy}{ed-c3}{énoncé : TD, p.~3}
D'après l'exemple de résonance, $y=-\frac t2\cos t+A\cos t+B\sin t$. $y(0)=A=0$ ; $y'(t)=-\frac12\cos t+\frac t2\sin t+B\cos t$,
$y'(0)=B-\frac12=0$. La solution est $y=\frac{\sin t-t\cos t}2$. Elle n'est pas bornée : $y(2k\pi)=-k\pi\to-\infty$.
\par\hfill\textcolor{BleuClair}{$\square$}\par
\partie{II}{Exercices classiques et méthodes incontournables}{}
\correction{9}{Seconds membres et paramètre}{ed-c2}{énoncé : TD, p.~4}
\q{(a)} Racine double $-2$ : $S_H=\{(A+Bt)e^{-2t}\}$. Pour $z''+4z'+4z=e^{it}$, $i$ n'est pas racine : $z=\frac{e^{it}}{(i+2)^2}=\frac{e^{it}}{3+4i}=\frac{(3-4i)e^{it}}{25}$,
d'où $y_p=\operatorname{Im}z=\frac{3\sin t-4\cos t}{25}$ et $y=(A+Bt)e^{-2t}+\frac{3\sin t-4\cos t}{25}$.

\q{(b)} Racine double $1$ : $S_H=\{(A+Bt)e^t\}$. Si $m\neq1$, $m$ n'est pas racine : $y_p=\frac{e^{mt}}{(m-1)^2}$. Si $m=1$, racine
double : $Q''=1$, $y_p=\frac{t^2}2e^t$.
\retenir{l'expression de $y_p$ change de forme exactement quand $m$ devient racine : discuter selon la multiplicité.}
\par\hfill\textcolor{BleuClair}{$\square$}\par
\correction{10}{Racine double et second membre $(t+1)e^{2t}$}{edn-3}{énoncé : TD, p.~4}
$r^2-4r+4=(r-2)^2$ : racine double $2$, $S_H=\{(A+Bt)e^{2t}\}$. Le second membre est $P(t)e^{2t}$ avec $2$ racine double :
on cherche $y=Q(t)e^{2t}$ et, comme $m^2+am+b=0$ et $2m+a=0$, il reste $Q''=t+1$ : $Q=\frac{t^3}6+\frac{t^2}2$ (on peut prendre
les constantes nulles). Solutions : $y=\left(\frac{t^3}6+\frac{t^2}2+A+Bt\right)e^{2t}$.
\retenir{racine double : $Q''=P$, il suffit d'intégrer deux fois.}
\par\hfill\textcolor{BleuClair}{$\square$}\par
\correction{11}{Superposition : second membre $\sh(2t)$}{edn-5}{énoncé : TD, p.~4}
Racines $1$ et $2$ : $S_H=\{Ae^t+Be^{2t}\}$. On écrit $\sh(2t)=\frac12e^{2t}-\frac12e^{-2t}$ et l'on superpose.
Pour $\frac12e^{2t}$ ($m=2$ racine simple) : $Q''+(2m+a)Q'+(m^2+am+b)Q=Q''+Q'=\frac12$, $Q=\frac t2$.
Pour $-\frac12e^{-2t}$ ($m=-2$ non racine) : $(4+6+2)Q=-\frac12$, $Q=-\frac1{24}$.
Solutions : $y=\frac t2e^{2t}-\frac1{24}e^{-2t}+Ae^t+Be^{2t}$.
\piege{$\sh(2t)$ n'est pas « de la forme » d'une seule exponentielle : on le décompose avant de chercher.}
\par\hfill\textcolor{BleuClair}{$\square$}\par
\correction{12}{Variation de la constante au premier ordre}{edn-17}{énoncé : TD, p.~4}
Forme normalisée : $y'+\frac t{1+t^2}y=\frac1{1+t^2}$ ; une primitive de $\frac t{1+t^2}$ est $\frac12\ln(1+t^2)$, d'où
$S_H=\left\{\frac K{\sqrt{1+t^2}}\right\}$. Aucune solution évidente : on pose $y=\frac{K(t)}{\sqrt{1+t^2}}$, et il reste
$\frac{K'(t)}{\sqrt{1+t^2}}=\frac1{1+t^2}$, soit $K'=\frac1{\sqrt{1+t^2}}$, $K=\argsh t+C$. Solutions :
\[y=\frac{\argsh t+C}{\sqrt{1+t^2}}=\frac{\ln\big(t+\sqrt{1+t^2}\big)+C}{\sqrt{1+t^2}}.\]
\retenir{en variation de la constante, les termes en $K$ se simplifient toujours : il ne reste que $K'\varphi=$ second membre.}
\par\hfill\textcolor{BleuClair}{$\square$}\par
\correction{13}{Une équation intégrale}{ed-integ}{énoncé : TD, p.~5}
\choix{l'inconnue est sous une intégrale de borne variable : on dérive (méthode~1.2 du livret).}
Si $f$ convient, $f(x)=1+x+\int_0^xtf(t)\,dt$ est dérivable, $f'(x)=1+xf(x)$, et $f(0)=1$. Réciproquement, une solution de ce
problème de Cauchy, intégrée de $0$ à $x$, vérifie l'équation de départ. On résout $y'-xy=1$ : l'homogène donne $Ke^{x^2/2}$ ;
la variation de la constante donne $\lambda'=e^{-x^2/2}$, sans primitive usuelle. On écrit donc la solution du problème de Cauchy
sous forme intégrale :
\[f(x)=e^{x^2/2}\left(1+\int_0^xe^{-u^2/2}\,du\right).\]
\retenir{une primitive sans expression usuelle s'écrit $\int_{x_0}^x$ ; la borne $x_0$ est choisie pour que la condition initiale soit simple.}
\par\hfill\textcolor{BleuClair}{$\square$}\par
\correction{14}{Résonance selon la pulsation}{edn-8}{énoncé : TD, p.~5}
$S_H=\{A\cos\omega t+B\sin\omega t\}$. On résout $z''+\omega^2z=e^{it}$.
\begin{itemize}
\item $\omega\neq1$ : $i$ n'est pas racine de $r^2+\omega^2$, $z=\frac{e^{it}}{\omega^2-1}$, d'où $y_p=\frac{\cos t}{\omega^2-1}$.
Toutes les solutions sont bornées.
\item $\omega=1$ : $i$ est racine simple, $z=Q(t)e^{it}$ avec $Q''+2iQ'=1$, $Q=-\frac i2t$ ; $y_p=\operatorname{Re}z=\frac t2\sin t$.
Aucune solution n'est bornée (le terme $\frac t2\sin t$ ne l'est pas).
\end{itemize}
\retenir{excitation à la pulsation propre : résonance, amplitude non bornée.}
\par\hfill\textcolor{BleuClair}{$\square$}\par
\partie{III}{Exercices de synthèse et approfondissement}{}
\correction{15}{Une équation fonctionnelle}{edn-18}{énoncé : TD, p.~6}
\choix{on se ramène à une équation différentielle en dérivant la relation.}
Soit $f$ une solution. Comme $x\mapsto f(-x)$ est dérivable, $f'$ l'est : $f$ est deux fois dérivable et
$f''(x)=-f'(-x)=-f(x)$. Donc $f=A\cos+B\sin$. Réciproquement, $f'(x)=-A\sin x+B\cos x$ et $f(-x)=A\cos x-B\sin x$ :
l'égalité pour tout $x$ équivaut à $B=A$. Les solutions sont les fonctions $x\mapsto A(\cos x+\sin x)$.
\piege{l'équation $f''=-f$ n'est qu'une condition \emph{nécessaire} : la réciproque élimine une constante.}
\par\hfill\textcolor{BleuClair}{$\square$}\par
\correction{16}{Développement limité d'une solution}{edn-7}{énoncé : TD, p.~6}
\choix{on n'a pas besoin de $f$ : l'équation fournit une relation entre les coefficients du développement.}
Sur $]-\infty,1[$, $f'=\frac{\sin2t+t^2-f}{2(t-1)}$ ; par récurrence, $f$ est de classe $\mathcal C^\infty$, donc admet un
développement $f(t)=\sum_{k=0}^5a_kt^k+o(t^5)$, et $f'(t)=\sum_{k=1}^5ka_kt^{k-1}+o(t^4)$. Le coefficient de $t^k$ dans
$2tf'-2f'+f$ est $(2k+1)a_k-2(k+1)a_{k+1}$, et $\sin2t+t^2=2t+t^2-\frac43t^3+o(t^4)$. Avec $a_0=f(0)=0$ :
\[a_0-2a_1=0,\quad3a_1-4a_2=2,\quad5a_2-6a_3=1,\quad7a_3-8a_4=-\tfrac43,\]
d'où $a_1=0$, $a_2=-\frac12$, $a_3=-\frac7{12}$, $a_4=-\frac{11}{32}$ :
\[f(t)=-\frac{t^2}2-\frac{7t^3}{12}-\frac{11t^4}{32}+o(t^4).\]
\par\hfill\textcolor{BleuClair}{$\square$}\par
\partie{IV}{Exercices difficiles et type concours}{}
\correction{17}{Solutions périodiques d'un oscillateur forcé}{edn-15}{énoncé : TD, p.~7}
\q{1.} $y_0(t)=\sin t\int_0^tf(u)\cos u\,du-\cos t\int_0^tf(u)\sin u\,du$. En dérivant (produit et intégrale de borne variable),
les termes $\sin t\,f(t)\cos t-\cos t\,f(t)\sin t$ s'annulent : $y_0'(t)=\cos t\int_0^tf\cos+\sin t\int_0^tf\sin$. En dérivant
encore : $y_0''=-y_0+f(t)(\cos^2t+\sin^2t)=-y_0+f$. Enfin $y_0(0)=y_0'(0)=0$.

\q{2.} Comme $f(u)\cos u$ est $2\pi$-périodique, $\int_0^{t+2\pi}f\cos=\int_0^tf\cos+\int_t^{t+2\pi}f\cos=\int_0^tf\cos+C$, et de même avec
$\sin$ ; en reportant dans l'expression de $y_0$, on obtient $y_0(t+2\pi)=y_0(t)+C\sin t-S\cos t$.

\q{3.} Les solutions sont $y=y_0+A\cos t+B\sin t$, et $y(t+2\pi)-y(t)=C\sin t-S\cos t$ pour toutes. Cette fonction est nulle si et
seulement si $C=S=0$ (prendre $t=0$ et $t=\frac\pi2$). Donc ou bien toutes les solutions sont $2\pi$-périodiques, ou bien aucune.

\q{4.} $f=\cos$ : $C=\pi\neq0$, aucune solution périodique (résonance : on retrouve $\frac t2\sin t$). $f(t)=\cos2t$ : $C=S=0$, toutes
les solutions sont $2\pi$-périodiques (par exemple $y_p=-\frac13\cos2t$).
\par\hfill\textcolor{BleuClair}{$\square$}\par
\end{document}
